\section{Conclusion}

We have constructed a gauge-theoretic representation of quantum states using $\mathbb{Z}_8$ gauge theory with topological defects. Our main findings are:

\begin{enumerate}
    \item \textbf{Theorem~\ref{thm:1d-polynomial}:} In one dimension, expectation values of local observables are computable in polynomial time for states generated by Clifford+T circuits with $O(\text{poly}(n))$ $T$ gates. The key insight is the factorization of the partition function into a product of $8\times8$ transfer matrices.
    
    \item \textbf{Theorem~\ref{thm:2d-hard}:} In two dimensions, the evaluation problem is \#P-hard due to exponential slice dimension. The slice dimension grows as $8^{\Omega(n)}$, making even writing down the transfer matrix exponential in $n$.
    
    \item \textbf{Defect-Complexity Conjecture:} For generic quantum states, abelian $\mathbb{Z}_8$ gauge theory requires an exponential number of defects to represent them exactly. We provided an information-theoretic lower bound and argued that random states require exponentially many defects.
    
    \item \textbf{Dimensional Threshold:} We identified a clear dimensional threshold: 1D chains admit polynomial-time evaluation, while 2D and higher-dimensional lattices are \#P-hard.
\end{enumerate}

\subsection{Summary of Contributions}

\begin{itemize}
    \item \textbf{Conceptual:} Unified language connecting amplituhedron, tensor networks, and quantum shadows via gauge theory.
    \item \textbf{Technical:} Explicit dictionary between qubits and gauge theory variables, including precise definitions of Gauss law constraints and defect charges.
    \item \textbf{Theoretical:} Rigorous proof of polynomial evaluation in 1D and \#P-hardness in 2D, with clear separation between proven results and conjectures.
\end{itemize}

\subsection{Outlook}

Future directions include:
\begin{itemize}
    \item Proving or disproving the Defect-Complexity Conjecture for circuit-prepared states.
    \item Constructing explicit non-abelian gauge theory models with polynomial evaluation.
    \item Exploring connections to quantum gravity (defects as particles, partition functions as spin-foam amplitudes).
    \item Developing practical classical algorithms for 1D quantum circuits based on the transfer matrix method.
\end{itemize}

This work opens a new avenue for understanding the boundary between classical and quantum computation through geometric representations. The dimensional threshold we identified may have implications for understanding when quantum advantage becomes possible.
